\par\noindent\textbf{(1)} 合光场 \(E=E_1+e^{i\phi}E_2\)，故平均强度 \(I(\phi)=I_1+I_2+2\operatorname{Re}[e^{i\phi}\Gamma_{12}]\)。该式由复振幅二阶矩展开得到，若源场相关就不能将平均强度简单相加。
\par\noindent\textbf{(2)} 定义 \(\gamma_{12}=\Gamma_{12}/\sqrt{I_1I_2}\)。相干矩阵 \(J=\begin{pmatrix}I_1&\Gamma_{12}\\\Gamma_{12}^*&I_2\end{pmatrix}\) 是协方差矩阵，必须半正定；由行列式非负得 \(I_1I_2-|\Gamma_{12}|^2\ge0\)，因此 \(|\gamma_{12}|\le1\)。
\par\noindent\textbf{(3)} 选择 \(\phi\) 与相关相位同向或反向，得到 \(I_{\max,\min}=I_1+I_2\pm2|\Gamma_{12}|\)。于是条纹可见度 \(V=(I_{\max}-I_{\min})/(I_{\max}+I_{\min})=2\sqrt{I_1I_2}|\gamma_{12}|/(I_1+I_2)\)。
\par\noindent\textbf{(4)} 若两个源完全相干且等强，\(V=1\)，暗纹可为零；若互相关为零，\(V=0\)，不再干涉。即使 \(|\gamma_{12}|=1\)，不等强的源也有 \(V<1\)，所以不能把低可见度全部归因于随机相位不相干。
